\section{Weighted graph deformations and a sweep}
\label{sec:graphs}

This section develops a deformation that changes the metric while
preserving a weighted volume density.  Its input is a smooth function
on a Riemannian manifold; its output is the metric induced on the
graph of that function.  We first compute the change of weighted
scalar curvature, then solve graph equations prescribing weighted mean
curvature to control this change.  A localized version of the construction will give an
estimate for one sweep through the coordinates of a root probability
map.

For a metric \(h\) and a smooth weight \(\phi\), write
\[
 \wscal(h,\phi)=\Scal_h-2\Delta_h\phi-|d\phi|_h^2,
 \qquad
 \diver_\phi V=e^{-\phi}\diver(e^\phi V),
 \qquad
 \Delta_\phi=\diver_\phi\nabla .
\]
Thus the weighted operators use \(e^\phi dV_h\).  The scalar quantity
is Perelman's \(R^m\), with his \(f=-\phi\)
\cite[Section~1.3]{Perelman2002}.  The estimates below concern its
behavior under particular graph deformations.  For antecedents of
the graph and vertical Jacobi identities, see
\cite[Proposition~4.1]{OpenAIScalarInessentiality2026}; we include the
calculation with the conventions needed here.

\subsection{The graph identity}

\begin{lemma}[Weighted graph identity]
\label{lem:graph-identity}
Let \((X,h)\) be a smooth Riemannian manifold, let \(\phi,f\) be
smooth real functions, and let \(N>0\).  Set
\[
 L=(1+N^2|df|_h^2)^{1/2},\qquad v=L^{-1},\qquad
 U=\frac{N\nabla_h f}{L},\qquad
 h'=h+N^2df\otimes df,\qquad \phi'=\phi-\log L.
\]
View \(h'\) as the induced metric on the graph
\(\Sigma=\{(x,Nf(x)):x\in X\}\) in
\((X\times\R,h+dz^2)\), and extend \(\phi\) constantly in the
\(z\)-direction.  Orient the graph by
\(\nu=(-U,v)\), put
\(\II(V,W)=\langle\bar\nabla_V\nu,W\rangle\), and let
\[
 H=\tr_\Sigma\II,\qquad
 D=H+\partial_\nu\phi=-\diver_\phi U .
\]
Here \(D\), and its derivative \(U(D)\), are regarded as functions on
\(X\) through the graph projection.  Then
\begin{equation}
 \wscal(h',\phi')-\wscal(h,\phi)
 =D^2-2U(D)+|\II|_{h'}^2+|d\log L|_{h'}^2.
 \label{eq:graph-identity}
\end{equation}
Moreover,
\begin{equation}
 e^{\phi'}dV_{h'}=e^\phi dV_h,\qquad
 N^2|df|_{h'}^2=1-v^2=|U|_h^2,\qquad
 |d\log L|_{h'}\le L|U|_h|\II|_{h'} .
 \label{eq:graph-basic-bounds}
\end{equation}
\end{lemma}

\begin{proof}
Write \(\bar h=h+dz^2\), and suppress pullback by the graph
projection.  The formula for the graph mean curvature gives
\(H=-\diver U\), while \(\partial_\nu\phi=-U(\phi)\); hence the stated
formula for \(D\) holds.  The rank-one determinant formula gives
\(dV_{h'}=L\,dV_h\), proving the first assertion in
\eqref{eq:graph-basic-bounds}.  The inverse rank-one formula gives its
second assertion.

Let
\[
 P=|\II|^2+
   \bigl(\Ric_{\bar h}-\bar\nabla^2\phi\bigr)(\nu,\nu).
\]
The Gauss equation and the decompositions of the Laplacian and
gradient of \(\phi\) give
\begin{align*}
 \Scal_{h'}&=\Scal_{\bar h}-2\Ric_{\bar h}(\nu,\nu)
                 +H^2-|\II|^2,\\
 \Delta_{\bar h}\phi
   &=\Delta_\Sigma(\phi|_\Sigma)
       +H\partial_\nu\phi+\bar\nabla^2\phi(\nu,\nu),\\
 |d(\phi|_\Sigma)|^2&=|d\phi|_{\bar h}^2-(\partial_\nu\phi)^2.
\end{align*}
Since \(\wscal(\bar h,\phi)=\wscal(h,\phi)\), their combination is
\begin{equation}
 \wscal(h',\phi|_\Sigma)
 =\wscal(h,\phi)+D^2-2P+|\II|^2.
 \label{eq:graph-gauss-weighted}
\end{equation}
Adding \(\log v\) to the restricted weight changes the left side by
\[
 -2\Delta_{\Sigma,\phi}\log v-|d\log v|^2
 =-\frac{2\Delta_{\Sigma,\phi}v}{v}+|d\log v|^2 .
\]

The tangential part of vertical translation is
\[
 T=\partial_z-v\nu=(vU,1-v^2).
\]
For a hypersurface variation with normal speed \(v\), the variation
of \(H\) is
\(-\Delta_\Sigma v-(|\II|^2+\Ric_{\bar h}(\nu,\nu))v\).
The additional variation of \(\partial_\nu\phi\) is
\[
 v\bar\nabla^2\phi(\nu,\nu)
       -\langle d(\phi|_\Sigma),dv\rangle_\Sigma .
\]
Including the tangential velocity \(T\), and using that vertical
translation preserves the metric, weight, and \(D\), therefore gives
\begin{equation}
 (\Delta_{\Sigma,\phi}+P)v=T(D).
 \label{eq:graph-translation}
\end{equation}
The horizontal part of \(T\) is \(vU\), so \(T(D)=vU(D)\).
Substitution in \eqref{eq:graph-gauss-weighted} proves
\eqref{eq:graph-identity}.  Finally, differentiating
\(v=\langle\nu,\partial_z\rangle\) tangentially gives
\(dv(V)=\II(V,T)\).  Since \(|T|^2=1-v^2=|U|^2\), it follows that
\(|d\log v|\le v^{-1}|U||\II|\), which is the last bound in
\eqref{eq:graph-basic-bounds}.
\end{proof}

\paragraph{A one-coordinate calculation.}
The square and differentiated terms in the graph identity suggest
an ideal target for \(ND\), the scaled weighted mean curvature.
Let \(F\) be one smooth input coordinate, restrict
to a region where \(f>0\), and fix \(\epsilon,a>0\).  Also let
\(h_{\mathrm{in}}\le h\) be a comparison metric.  In the later sweep,
the incoming energy is measured in \(h_{\mathrm{in}}\), while \(h\)
already includes the earlier graphs.

Put \(A=|U|_h\).  Where \(A>0\), write \(U=Ae\) with \(|e|_h=1\).
The longitudinal unit tangent to the graph is
\(\tau=(ve,A)\).  Define
\[
 x=-\II(\tau,\tau),\qquad X_{\parallel}=Nx,\qquad t=A^2/f,\qquad
 B=1-|e|_{h_{\mathrm{in}}}^2\in[0,1].
\]
At \(A=0\), choose any \(h\)-unit vector \(e\), put \(\tau=(e,0)\),
and use the same definitions.  With \(\tau^\flat\) denoting the dual
covector in the graph metric \(h'\), set
\[
 \II^{\mathrm{rem}}=\II+x\,\tau^\flat\otimes\tau^\flat .
\]
The component of \(\II^{\mathrm{rem}}\) in the
\(\tau^\flat\otimes\tau^\flat\) direction vanishes, so
\begin{equation}
 \begin{gathered}
 NU(f)=LA^2,\qquad
 |\II|_{h'}^2=x^2+|\II^{\mathrm{rem}}|_{h'}^2,\qquad
 \tr_{h'}\II^{\mathrm{rem}}=H+x,\\
 x=e(A)\quad\text{where }A>0.
 \end{gathered}
 \label{eq:graph-longitudinal}
\end{equation}
For the last identity, use \(\nu=(-Ae,v)\) and \(A^2+v^2=1\) to
compute
\[
 \II(\tau,\tau)=-v^2e(A)+Av\,e(v)=-e(A).
\]
The first identity in \eqref{eq:graph-longitudinal} follows directly
from the definitions of \(U,L,A\).

Define the ideal target for \(ND\) by
\begin{equation}
 Z=\epsilon a(F^2/f-f)+A^2/(2f).
 \label{eq:graph-ideal-target}
\end{equation}
Since \(NU(A^2)=2A^2X_{\parallel}\), with both sides zero at \(A=0\),
differentiation gives
\begin{align}
 -2NU(Z)
 ={}&-4\epsilon a(F/f)NAe(F)
       +2\epsilon a(F^2/f^2+1)LA^2 \notag\\
    &-2X_{\parallel}t+Lt^2.
 \label{eq:graph-ideal-derivative}
\end{align}
The first term couples the prescribed \(F^2/f\) to the incoming
energy.  The \(A^2/(2f)\) term completes the longitudinal square:
\begin{equation}
 \begin{gathered}
 X_{\parallel}^2-2X_{\parallel}t+Lt^2
   =(X_{\parallel}-t)^2+(L-1)t^2,\\
 \begin{aligned}
 &\epsilon N^2|dF|_{h_{\mathrm{in}}}^2
       -4\epsilon a(F/f)NAe(F)\\
 &\quad=\epsilon\left|N\,dF
       -2a(F/f)A\,h_{\mathrm{in}}(e,\cdot)\right|_{h_{\mathrm{in}}}^2
       -4\epsilon a^2(F/f)^2A^2(1-B)\\
 &\quad\ge-4\epsilon a^2(F/f)^2A^2(1-B).
 \end{aligned}
 \end{gathered}
 \label{eq:graph-ideal-completions}
\end{equation}
Consequently
\begin{align}
 &Z^2-2NU(Z)+N^2|\II|_{h'}^2+N^2|d\log L|_{h'}^2
                 +\epsilon N^2|dF|_{h_{\mathrm{in}}}^2 \notag\\
 \ge{}&Z^2+(X_{\parallel}-t)^2+N^2|\II^{\mathrm{rem}}|_{h'}^2
                 +N^2|d\log L|_{h'}^2 \notag\\
 &+\epsilon A^2\bigl[(2a-4a^2)F^2/f^2+2a\bigr]
       +4\epsilon a^2(F^2/f^2)A^2B \notag\\
 &+(L-1)\bigl[2\epsilon a(F^2/f^2+1)A^2+t^2\bigr].
 \label{eq:graph-ideal-bound}
\end{align}
At \(A=0\), \(t=0\) and
\((X_{\parallel}-t)^2+N^2|\II^{\mathrm{rem}}|^2=N^2|\II|^2\); every other
direction-dependent term vanishes.  Thus the inequality is independent
of the arbitrary direction at such points.

This calculation is algebraic and assumes no equation for \(D\).
The complete equation below realizes \(ND=Z-\beta f\) up to a positive
floor, and localization introduces controlled cutoff errors.  A sweep
will also have to account for normalization of the output map and the
interaction of the ordered graph metrics.

\subsection{Prescribed graphs on complete leaves}

We specify the uniformity required for the analytic construction.
A manifold has \emph{bounded geometry} here if it has coordinate
charts of a common positive size whose smaller concentric domains
cover the manifold, with a common positive buffer to each chart
boundary, and with uniform bounds on the metric, its inverse, and
all their coordinate derivatives.  These bounds, as well as the
chart size, may depend on the instance.  A function is uniformly
smooth if all its coordinate derivatives, including its value, are
uniformly bounded on these charts.  For a weight \(\phi\), only
positive-order derivatives are required to be uniformly bounded.
The \emph{uniform Hölder space} below means the Banach space
obtained by taking the supremum of the corresponding Hölder norms in
these charts, after a fixed shrinkage within the buffers.

The use of complete prescribed graphs has precedents in
\cite[Theorem~5.2]{OpenAIScalarInessentiality2026}.  The
following statement records the precise uniform and measurable
conclusions used here and in the hypersurface argument.  Its positive
branch realizes the designed prescription with a positive floor and
will be localized for the sweep.  Its tangent branch supplies the value
approximation and the intrinsic gradient bound used in
Section~\ref{sec:degree}.

\begin{lemma}[Prescribed graphs on a complete manifold]
\label{lem:graph-prescription}
Let \((X,h)\) be a complete smooth Riemannian manifold without
boundary, with bounded geometry as above, and suppose the
positive-order derivatives of \(\phi\) are uniformly bounded.  For
a graph of height \(z\), set
\[
 L_z=(1+|dz|_h^2)^{1/2},\quad v_z=L_z^{-1},\quad
 U_z=\frac{\nabla_hz}{L_z},\quad
 D_z=-\diver_\phi U_z .
\]
Consider either of the following equations \(D_z=p(x,z,v_z)\).

\begin{enumerate}
\renewcommand{\labelenumi}{\textup{(\roman{enumi})}}
\item Let \(N,\epsilon,a>0\), \(0<\delta\le1\), and \(H_*\ge0\).
  Suppose \(F_i,\beta\) are uniformly smooth,
  \(|F_i|\le1\), and \(0\le\beta\le H_*\).  On the region \(z>0\),
  prescribe
  \begin{equation}
   p(x,z,v)=
   \frac{\epsilon a(F_i(x)^2+\delta^2)+(1-v^2)/2}{z}
       -(\epsilon a+\beta(x))\frac{z}{N^2}.
   \label{eq:graph-prescription-positive}
  \end{equation}
\item Let \(e,w>0\) be constants.  Suppose \(B,f_0\) are uniformly
  smooth and \(B\ge B_{\min}>0\).  On the region
  \(|(f_0(x)-ez)/w|<\pi/2\), prescribe
  \begin{equation}
   p(x,z,v)=B(x)\tan\bigl((f_0(x)-ez)/w\bigr).
   \label{eq:graph-prescription-tangent}
  \end{equation}
\end{enumerate}

In either case there is a smooth solution of bounded height with
uniform bounds on all its derivatives, and it lies uniformly inside
the stated open prescription region.  Its angle satisfies
\(v_z\ge v_0>0\).  In Case \textup{(i)} it satisfies
\begin{equation}
 c_{\min}:=\delta\sqrt{\frac{\epsilon a}{\epsilon a+H_*}}
       \ \le\ f:=z/N\ \le 2.
 \label{eq:graph-height-positive}
\end{equation}
In Case \textup{(ii)} there is \(\eta>0\) such that
\[
 \left|\frac{f_0-ez}{w}\right|\le\frac{\pi}{2}-\eta;
 \qquad
 |ez-f_0|<\frac{\pi w}{2},\quad
 |z|\le\frac{\|f_0\|_\infty+\pi w/2}{e}.
\]
The constants \(v_0,\eta\) and the derivative bounds may depend on
all the indicated parameters and common bounds of the data.
The solution is unique among bounded-height solutions whose
\(C^{2,\alpha}\) norms in all buffered charts are uniformly bounded
for some \(0<\alpha<1\), and whose values have a uniform margin inside
the prescription region.  Each such solution may have its own
exponent, bound, and margin.

The graph metric \(h_z=h+dz\otimes dz\) is again complete and has
bounded geometry, and \(\phi_z=\phi-\log L_z\) has uniform
positive-order derivatives.  In particular,
\[
 e^{\phi_z}dV_{h_z}=e^\phi dV_h,\qquad
 |dz|_{h_z}^2=1-v_z^2<1.
\]
For the family assertion, let \(\mathcal T\) be a standard Borel
parameter space indexing complete problems of a fixed dimension.
Assume a common chart radius and buffers, common bounds for the
metric, its inverse, and all their chart derivatives, and common
bounds for the positive-order derivatives of \(\phi\).  Require
common bounds, including order zero, for all chart derivatives of
\(F_i,\beta\) in Case \textup{(i)} or \(B,f_0\) in Case
\textup{(ii)}.  In Case \textup{(i)}, let \(N,\epsilon,a,\delta\)
range in fixed compact subsets of their permitted positive ranges,
and require a common finite upper bound on \(H_*\).  In Case
\textup{(ii)}, let \(e,w\) range in fixed compact subsets of
\((0,\infty)\), and require \(B\ge B_{\min}\) for one common
\(B_{\min}>0\).  Suppose also that the data and chart overlaps are
Borel and that the charts accessible along each leaf admit a
countable Borel enumeration.  Then the preceding estimates for the
solutions and their graph metrics and weights hold with constants common to the family,
including a common positive height margin in Case \textup{(i)}
and a common margin from the poles in Case \textup{(ii)}.  The
solutions are Borel in the chart representations and commute with
any equivariant identifications of the complete problems.
\end{lemma}

\begin{proof}
We give existence and the estimates together, since the estimates
also give uniqueness and measurable dependence.

\emph{Height barriers and comparison.}
In Case \textup{(i)}, replace \(F_i,\beta\) by \(tF_i,t\beta\),
\(0\le t\le1\).  At \(t=0\) the constant \(f=\delta\) is a
solution.  For every \(t\), the constant graphs \(f=c_{\min}\) and
\(f=2\) are a lower and an upper barrier, respectively.  Indeed their
weighted mean curvature is zero, and at angle \(v=1\) the right side
of \eqref{eq:graph-prescription-positive} is nonnegative at
\(f=c_{\min}\) and strictly negative at \(f=2\).  In Case
\textup{(ii)}, replace \(f_0\) by \(tf_0\), starting with \(z=0\).
For \(\tau=\pi/2-\eta\), the graphs
\[
 z_-=\frac{tf_0-w\tau}{e},\qquad
 z_+=\frac{tf_0+w\tau}{e}
\]
have arguments \(+\tau\) and \(-\tau\), respectively.  Their
weighted mean curvatures have a common bound independent of
\(\eta\), because their derivatives do not involve \(\tau\).
Choose \(\eta>0\) so small that \(B_{\min}\tan\tau\) exceeds this
bound.  These are then lower and upper barriers for every \(t\).
They give the asserted height and pole margins.

Both comparisons use strict decrease of \(p\) in its height
argument, with the other arguments held fixed:
\begin{align}
 p_z&=-\frac{\epsilon a(F_i^2+\delta^2)+(1-v^2)/2}{z^2}
        -\frac{\epsilon a+\beta}{N^2}
       \le-\frac{\epsilon a}{N^2}
       &&\text{in Case \textup{(i)}}, \notag\\
 p_z&=-\frac{eB}{w}\sec^2\bigl((f_0-ez)/w\bigr)
       \le-\frac{eB_{\min}}w
       &&\text{in Case \textup{(ii)}}.
 \label{eq:graph-height-monotonicity}
\end{align}
At an interior positive maximum of the difference between a
solution and an upper barrier, their gradients, hence their angles,
agree.  The negative elliptic principal part of \(D_z\) gives a
nonnegative difference of mean curvatures there, whereas strict
height decrease gives a negative difference of prescriptions.
This is impossible.  The lower comparison and uniqueness follow in
the same way.  On a noncompact manifold the maximum can be replaced
by a supremum: take uniform charts centered at an extremizing
sequence, subtract the value of \(\phi\) at each center, and pass to
a local limit.  For any individual solution in the stated uniqueness
class, its uniform \(C^{2,\alpha}\) bound and the coefficient bounds
give such a limit, which attains the extremum at the center of an
interior chart.  This
justifies the same comparison without compactness.

\emph{The angle bound.}
For a solution, \eqref{eq:graph-translation} and the chain rule give
\begin{equation}
 (\Delta_{\Sigma,\phi}+P)v_z
   =v_z\langle p_x,U_z\rangle
       +p_z(1-v_z^2)+p_vT(v_z).
 \label{eq:graph-angle-equation}
\end{equation}
Here \(p_x\) differentiates with \(z,v\) fixed, and we used
\(T(z)=1-v_z^2\).  The ambient curvature and weight bounds imply
\(P\ge-C\), since the remaining term \(|\II|^2\) is nonnegative.
At an infimum of \(v_z\), taken by the same recentering procedure,
\(\Delta_{\Sigma,\phi}v_z\ge0\) and \(T(v_z)=0\).  Consequently, if
\(v_z\le1/2\) there,
\begin{equation}
 (1-v_z^2)|p_z|\le v_z(C+|p_x|).
 \label{eq:graph-angle-infimum}
\end{equation}
The height margins just proved bound \(p_x\) uniformly, and
\eqref{eq:graph-height-monotonicity} bounds \(|p_z|\) below by a
positive constant.  Since \(1-v_z^2\ge3/4\) in this case,
\eqref{eq:graph-angle-infimum} gives a uniform positive lower bound
for \(v_z\).  In using recentered limits at this stage, a solution
is first in an individual uniform \(C^{2,\alpha}\) class.  Ordinary
local bootstrapping gives its individual higher regularity before
the common estimates are deduced, so the differentiations and
limits in this argument are justified.

\emph{Higher estimates.}
In a uniform chart let
\(\omega=e^\phi\sqrt{\det h}\), and divide \(\omega\) by its value
at the chart center.  The resulting density has a common upper
bound and a positive lower bound, and all its derivatives are
bounded.  The equation is of the form
\begin{equation}
 -\partial_j A^j(x,\partial z)=F(x,z,\partial z).
 \label{eq:graph-divergence-equation}
\end{equation}
For example \(A^j=\omega h^{jk}z_k/L_z\) after the normalization.
On the established height and gradient ranges the flux and the
right side are smooth with uniformly bounded derivatives, and
\(\partial A^j/\partial z_k\) is uniformly elliptic.  Differentiating
\eqref{eq:graph-divergence-equation} gives linear divergence
equations for \(\partial_kz\).  Their leading coefficients are
bounded and uniformly elliptic.  Gradient dependence of \(F\)
produces a bounded first-order coefficient; the remaining terms
are a bounded function plus the divergence of a bounded vector
field.  The term \(F_z\partial_kz\) is included in the bounded
source, since the gradient is already bounded.

The interior De Giorgi--Nash estimate for these equations, with
bounded lower-order coefficients and source data in \(L^{s/2}\)
and \(L^s\) for \(s>\dim X\), gives a common Hölder bound for
\(\partial z\); see \cite[Theorem~8.24]{GilbargTrudinger}.
The coefficients in the nondivergence form of the original
equation are now Hölder continuous.  Interior Schauder estimates
give a uniform \(C^{2,\alpha}\) bound for some \(\alpha>0\), after
shrinking within the chart buffers.  Differentiating and applying
Schauder estimates again yields uniform bounds for every
derivative.  All these estimates are local with common buffers, so
their constants are uniform throughout \(X\).

\emph{Continuation.}
Work in uniform \(C^{2,\alpha}\) spaces on the open set of functions
with a uniform height margin.  The equation \(D_z-p(x,z,v_z)=0\)
defines a smooth map to uniform \(C^{0,\alpha}\).  At a solution
its linearization has negative uniformly elliptic principal part,
uniform Hölder coefficients, and zeroth-order coefficient
\(-p_z\) bounded below by a positive constant.  This linear
operator is invertible on the uniform spaces.  To see
surjectivity, exhaust a noncompact component by smooth relatively
compact domains and solve the linear Dirichlet problem with zero
boundary data on each domain.  The maximum principle bounds the
supremum of the solutions by the supremum of the right side
divided by the positive zeroth-order lower bound.  Interior
Schauder estimates then give the same \(C^{2,\alpha}\) bound in
every buffered chart once it lies in the exhaustion.  A diagonal
local limit solves the equation on the component with the uniform
bound.  On a compact component the corresponding closed-manifold
elliptic theory gives the same conclusion.  Uniqueness follows
from the maximum principle, using recentered limits on noncompact
components.  The estimate also makes the inverse bounded.

The inverse function theorem gives openness of the set of
homotopy parameters with a solution.  The a priori height, angle,
and higher bounds are common over the homotopy.  For a sequence
of parameters converging to a limit, diagonal smooth convergence
on a countable chart cover therefore gives a solution at that
limit with the same bounds and height margin.  This gives
closedness.  Each homotopy starts from the constant solution
specified above, so existence follows.  Comparison gives the
claimed uniqueness.

The graph estimates now follow from
\eqref{eq:graph-basic-bounds}.  In particular \(h_z\ge h\), and
the gradient bound also gives a uniform upper comparison in the
given charts.  Thus \(h_z\) is complete and has uniform buffered
charts with bounds on all derivatives.  The corresponding
derivative bounds on \(\log L_z\) give those for \(\phi_z\).
This also permits successive applications of the lemma.

\emph{Borel dependence.}
For completeness, fix a starting plaque and enumerate the charts
accessible along its leaf by words in the Borel overlap maps.
Represent a solution by its functions in this countable family
of charts, in a countable product of smooth function spaces.
Compatibility on overlaps, the equation, the height margin, and
the common derivative bounds are Borel conditions: each can be
tested on countably many compact subcharts and dense sets, where
the relevant functions and their derivatives are continuous.
The chart data and changes of coordinates are Borel by
hypothesis.  The resulting Borel relation between data and local
representatives has a singleton fiber over each set of data,
by existence and uniqueness.  The Borel injection theorem
\cite[Corollary~15.2]{Kechris1995},
applied to its projection to the data, makes the inverse
projection Borel.  This proves measurable dependence.  The
coding refers to a full prescribed problem on the same complete
leaf.  It applies independently of the transverse mass of the
chart family.  Finally, applying an equivariant identification
to a solution produces a solution of the identified problem, so
uniqueness proves equivariance.
\end{proof}

\subsection{Localization of the positive prescription}

The positive solution in Lemma~\ref{lem:graph-prescription} is
defined on the complete manifold and need not have compact support.
For a probability coordinate we need a graph supported near that
coordinate.  We obtain one by forcing the complete solution close
to a small constant on an outer region, and cutting it off there.
The following lemma states the quantitative features of this
procedure that enter the sweep.

\begin{lemma}[Localization of the positive graph]
\label{lem:graph-localization}
Let \(X,h,\phi\) satisfy the hypotheses of
Lemma~\ref{lem:graph-prescription}, with \(X\) connected, and let
\(h_{\mathrm{ref}}\le h\) be a complete second metric.  Require the
coefficients of \(h_{\mathrm{ref}}\), its inverse, and all their
coordinate derivatives to have uniform bounds in the same buffered
charts used for \(h\).
Suppose \(F_i\) is uniformly smooth, \(|F_i|\le1\), with compact
support \(K\), and let \(\alpha\ge0\) be a uniformly smooth
function.  Fix
\[
 \epsilon>0,\qquad N^2=l\epsilon,\qquad
 a=\tfrac12+\theta,\qquad 0<l\le1,\qquad l\le\theta\le\tfrac1{16}.
\]
There is a radius \(R_{\mathrm{loc}}\), depending only on the
dimension and a fixed sufficiently small numerical constant
\(c_0>0\), with the following property.  For every sufficiently
large \(H_*\), one can choose smooth functions
\(\beta_s\ge0\) and \(0\le\chi\le1\) such that, with
\(\beta=\beta_s+\alpha\),
\begin{equation}
 0\le\beta\le H_*,
 \qquad
 N|d\beta_s|_h
       \le c_0\bigl(\theta\epsilon^{3/2}+\beta_s^{3/2}\bigr).
 \label{eq:graph-spatial-cutoff}
\end{equation}
The function \(\beta_s\) vanishes on a neighborhood of \(K\).
The function \(\chi\) is one there, is supported in
\(\{\dist_{h_{\mathrm{ref}}}(x,K)<R_{\mathrm{loc}}\}\), and changes
from one to zero only within a region where \(F_i=0\) and
\(\beta=H_*\).  This region and the support of \(d\chi\) have fixed
positive buffers in \(h_{\mathrm{ref}}\), independent of \(H_*\).
The derivatives of \(\chi\) are bounded independently of \(H_*\),
\(\chi\) is flat on its zero set, and \(|d\chi|_h^2\le C\chi\).

Set \(\delta=H_*^{-2}\), and let \(\widetilde f=z/N\) be the
complete solution of \eqref{eq:graph-prescription-positive} for
this \(\beta,\delta\).  Define \(f=\chi\widetilde f\) and form the
graph of height \(Nf\).  Denote its quantities by \(L,U,D,\II\),
put \(A=|U|_h\), and set, where \(f>0\),
\begin{equation}
 ND^0=\epsilon a(F_i^2/f-f)+A^2/(2f).
 \label{eq:graph-ideal-prescription}
\end{equation}
The expressions \(A^2/f\) and \(ND^0\) are extended by zero
where \(f=0\).  The bump \(f\) is smooth,
\(0\le f\le2\), is positive on \(K\), and is flat on its zero set.
On the region \(\chi=1\) it satisfies
\begin{equation}
 ND=ND^0-\beta f+\epsilon a\delta^2/f .
 \label{eq:graph-core-prescription}
\end{equation}
On the region where the cutoff changes, the following quantities
tend uniformly to zero as \(H_*\to\infty\), with the preceding
data and parameters fixed:
\begin{equation}
 \begin{gathered}
 N^2\bigl|\wscal(h+N^2df\otimes df,\phi-\log L)
                      -\wscal(h,\phi)\bigr|,\quad
 A^2/f,\quad |ND^0|,\quad \beta f,\\
 |ND|,\quad N|\II|,\quad N|d\log L|_{h+N^2df\otimes df}.
 \end{gathered}
 \label{eq:graph-cutoff-small}
\end{equation}
In particular, on all of \(X\) the function
\(\mathrm{err}:=ND-(ND^0-\beta f)\) tends uniformly to zero.
All data retain the corresponding uniform bounds for each fixed
choice of \(H_*\), with the positive-order convention for weights.
For the family assertion, let a standard Borel parameter space index
such problems of a fixed dimension.  Require common sizes and buffers
for these simultaneous charts, and common bounds there for
\(h,h^{-1},h_{\mathrm{ref}},h_{\mathrm{ref}}^{-1}\) and all their
coordinate derivatives.  Require common bounds for all derivatives,
including values, of \(F_i,\alpha\), and for the positive-order
derivatives of \(\phi\).  Let \((\epsilon,l,\theta)\) range in a fixed
compact subset of the permitted parameter set.  Require the data,
including \(h_{\mathrm{ref}}\), to
be Borel in the charts and preserved by the equivariant identifications
under consideration.  Require also the Borel chart overlaps and
countable Borel enumerations of accessible charts in
Lemma~\ref{lem:graph-prescription}.  Then the choices are Borel and
commute with those identifications; their bounds for each fixed
\(H_*\), and the limits above as \(H_*\to\infty\), are uniform over
the family.  When \(K\) is empty one may take \(f=0\) and omit the
operation.
\end{lemma}

\begin{proof}
We first construct the outer coefficient while keeping its
transition width bounded.  Starting at a fixed positive distance
from \(K\), a scalar profile \(b\) can increase from zero to
\(H_*\) at speed at most
\[
 \frac{c_0}{C N}\bigl(\theta\epsilon^{3/2}+b^{3/2}\bigr),
\]
where the fixed factor \(C\) leaves room for the estimates below.
Choose the speed comparable to this expression during the
increase, with smooth flattening at either end.  The width needed
is at most
\begin{equation}
 C N\int_0^\infty\frac{dt}
                  {\theta\epsilon^{3/2}+t^{3/2}}
 =C\frac{N}{\theta^{1/3}\sqrt\epsilon}
 =C\frac{\sqrt l}{\theta^{1/3}}\le C .
 \label{eq:graph-cutoff-width}
\end{equation}
For example one can use the primitive of the reciprocal of a smooth
positive function comparable to
\(\theta\epsilon^{3/2}+t^{3/2}\) as the new profile coordinate,
and compose with fixed smooth flattenings.  The last inequality
uses \(l\le\theta\le1/16\).

One may use this profile in a smooth approximation of
\(\dist_{h_{\mathrm{ref}}}(x,K)\) with gradient at most \(2\)
and a uniformly small error.  Here is a construction with the
needed uniformity.  Shrinking the controlled charts gives normal
charts at a fixed positive radius, by the uniform geodesic
equation estimates.  Average the distance function in such a
chart centered at \(x\) against a fixed smooth radial kernel in
\(T_xX\), at a sufficiently small radius.  Written as an integral
against a kernel in \((x,y)\), this is a smooth function of \(x\).
The error is at most the averaging radius.  To bound the first
derivative, move the center and parallel transport the averaging
vectors.  At a sufficiently small radius the displacement of
the exponential images is at most twice the displacement of the
center; the distance function is one-Lipschitz.  The derivative
bound follows.  Higher derivatives have uniform bounds from the
smooth local kernels: in differentiating the normalized kernel
one may subtract the distance at the center, leaving a uniformly
bounded oscillation on its support.  These derivative bounds hold in
the common buffered charts, since the \(h_{\mathrm{ref}}\) exponential
maps and kernels have uniformly bounded derivatives there.  The inverse metric
inequality \(h^{-1}\le h_{\mathrm{ref}}^{-1}\) then bounds the
gradient of the composed profile in \(h\).

Choose this raw profile \(b\) to vanish near \(K\), to equal
\(H_*\) past the width in \eqref{eq:graph-cutoff-width}, and to
leave fixed outer buffers.  To accommodate \(\alpha\), take a
smooth nondecreasing saturation
\(\sigma_{H_*}:[0,\infty)\to[0,H_*]\) that equals the identity
up to \(H_*/3\), equals \(H_*\) from \(H_*\) on, and has a
derivative bounded by a numerical constant.  Put
\[
 \beta=\sigma_{H_*}(b+\alpha),\qquad
 \beta_s=\sigma_{H_*}(b+\alpha)-\alpha .
\]
Take \(H_*\) much larger than the value and derivative bounds
of the already fixed \(\alpha\).  Then \(\beta_s\ge0\), and it
equals \(b\) near \(K\).  Wherever the saturation differs from
the identity, both \(\beta_s\) and the scale of \(b\) are
comparable to \(H_*\).  The formula
\[
 d\beta_s=\sigma_{H_*}'(b+\alpha)\,db
                  +\bigl(\sigma_{H_*}'(b+\alpha)-1\bigr)d\alpha
\]
therefore proves \eqref{eq:graph-spatial-cutoff}: choose the
initial fixed speed factor large enough to handle the first
term, and then choose \(H_*\) large enough for the fixed
\(d\alpha\) term to be a small fraction of
\(H_*^{3/2}/N\).  Where saturation is the identity this is just
the raw profile bound.  Where \(b=H_*\), the total \(\beta\)
equals \(H_*\).  A cutoff supported beyond fixed additional
buffers can consequently be chosen once at a radius
\(R_{\mathrm{loc}}\) as in the statement.  Take it to be the
square of a smooth nonnegative cutoff that is flat at its outer
end.  This gives \(|d\chi|^2\le C\chi\) and the other asserted
cutoff properties.  If some outer regions are empty, the
corresponding assertions are vacuous.

We next prove that the complete solution can be cut on this
plateau.  All constants in the rest of the proof may use the
fixed data and parameters, but not \(H_*\).  Write
\[
 c_H=H_*^{-2}\sqrt{\frac{\epsilon a}{\epsilon a+H_*}}
       \asymp H_*^{-5/2}.
\]
By \eqref{eq:graph-height-positive},
\(c_H\le\widetilde f\le2\).  Cover the cutoff region by inner chart
boxes, each with a fixed succession of larger buffered boxes whose
closures lie in the plateau \(F_i=0,\ \beta=H_*\).  Their sizes may
use the current \(h\)-chart bounds inside the fixed
\(h_{\mathrm{ref}}\) buffers; all these bounds were fixed before
\(H_*\), so the boxes and their buffers have sizes independent of
\(H_*\).  Our target is that, for every integer \(j\ge0\) and every
\(M>0\), there is \(C_{j,M}\) such that, for all sufficiently large
\(H_*\),
\begin{equation}
 \|\widetilde f-c_H\|_{C^j}
       \le C_{j,M}H_*^{-M},
 \label{eq:graph-plateau-jets}
\end{equation}
on the inner boxes, with norms in these fixed coordinates.
The constants use only the indicated data and parameters fixed
before \(H_*\).  The family assertion of
Lemma~\ref{lem:graph-prescription} is applied separately for each
fixed \(H_*\); the estimates that follow establish this
superpolynomial decay uniformly as \(H_*\) grows.

There is first a global angle bound
\begin{equation}
 v\ge\frac{c_1}{1+\sqrt{H_*}}
 \label{eq:graph-global-plateau-angle}
\end{equation}
for the uncut graph.  In fact, at a recentered infimum with
\(v\le1/2\), the formulas for \(p_x,p_z\) and
\eqref{eq:graph-spatial-cutoff} give
\[
 |p_z|\ge c_2(1+\beta+\widetilde f^{-2}),\qquad
 |p_x|\le C\bigl(\widetilde f^{-1}
                    +\widetilde f(1+\beta^{3/2})\bigr).
\]
For the first bound use \(1-v^2\ge3/4\) in \(p_z\); for the
second use the fixed bounds on \(F_i,\alpha\) and their
derivatives.  Since \(\widetilde f\le2\) and \(\beta\le H_*\),
the ratio of the right side of the second estimate, plus a
constant, to the right side of the first is at most
\(C(1+\sqrt{H_*})\).  Equation
\eqref{eq:graph-angle-infimum} proves
\eqref{eq:graph-global-plateau-angle}.  This recentering is done
for each fixed \(H_*\), using the individual regularity supplied
by Lemma~\ref{lem:graph-prescription}.

On each of the outer plateau boxes the equation is
\begin{equation}
 ND=\frac{\epsilon a\delta^2}{\widetilde f}
       +\frac{1-v^2}{2\widetilde f}
       -(\epsilon a+H_*)\widetilde f,
 \qquad
 (\epsilon a+H_*)c_H^2=\epsilon a\delta^2 .
 \label{eq:graph-plateau-equation}
\end{equation}
There is an exponential value estimate on an inner box:
\begin{equation}
 0\le\widetilde f-c_H\le C\exp(-c_3\sqrt{H_*}).
 \label{eq:graph-plateau-value}
\end{equation}
To prove it, let \(d_{\mathrm{face}}\) be ordinary coordinate
distance to a face of the outer box and put
\[
 w_1=\sum_{\mathrm{faces}}3
          \exp(-\gamma\sqrt{H_*}\,d_{\mathrm{face}}).
\]
On the boundary, \(c_H+w_1\ge3\ge\widetilde f\).
Its first and second coordinate derivatives are bounded by
\(C\gamma\sqrt{H_*}w_1\) and \(C\gamma^2H_*w_1\),
respectively.  Direct differentiation of the weighted graph
divergence therefore bounds its mean curvature from below by
\[
 -C(\gamma^2H_*+\gamma\sqrt{H_*})w_1 .
\]
For this graph, the height terms in its prescription satisfy
\[
 \frac{\epsilon a\delta^2}{N(c_H+w_1)}
     -\frac{(\epsilon a+H_*)(c_H+w_1)}N
 =-\frac{\epsilon a+H_*}{N}w_1
             \left(1+\frac{c_H}{c_H+w_1}\right).
\]
The remaining angle term is at most
\(C\gamma^2H_*w_1\), because
\((1-v^2)/(c_H+w_1)\le
 N^2|dw_1|^2/(c_H+w_1)\le C\gamma^2H_*w_1\).
Choosing \(\gamma>0\) sufficiently small and then \(H_*\)
sufficiently large makes \(c_H+w_1\) an upper barrier.
Comparison with this barrier and the lower constant \(c_H\)
proves \eqref{eq:graph-plateau-value} on any inner box a fixed
distance from the faces.

On a smaller box the angle is bounded below independently of
\(H_*\).  To see this, choose a smooth chart cutoff
\(0\le\xi\le1\), supported in a region where
\eqref{eq:graph-plateau-value} holds and equal to one on the
smaller box.  Minimize \(v/\xi\) over \(\{\xi>0\}\).
The global lower bound \eqref{eq:graph-global-plateau-angle}
makes a minimum with value below \(1/2\), if one exists,
interior, and at that minimum
\[
 v<\xi/2,\qquad \xi\ge c/\sqrt{H_*},\qquad
 \nabla_\Sigma v=(v/\xi)\nabla_\Sigma\xi,\qquad
 \Delta_{\Sigma,\phi}v\ge(v/\xi)\Delta_{\Sigma,\phi}\xi.
\]
The formula for the restriction of an ambient function to the
graph gives
\(\Delta_{\Sigma,\phi}\xi\ge-C(1+|\II|)\).  The positive
\(|\II|^2\) term in \(P\), and completion of its square against
\(C|\II|/\xi\), show that the left side of
\eqref{eq:graph-angle-equation} is at least \(-CH_*v\).
On this plateau \(p_x=0\).  Since \(v<1/2\), the \(p_z\) term
on its right side is at most \(-c/\widetilde f^2\).
Also \(p_v=-v/(N\widetilde f)\) and
\(|T(v)|\le Cv/\xi\), so that right side is at most
\[
 -\frac{c}{\widetilde f^2}
       +\frac{C\sqrt{H_*}\,v}{\widetilde f}.
\]
Combining the two bounds and multiplying by \(\widetilde f^2\)
would give
\[
 c\le CH_*\widetilde f^2+
                  C\sqrt{H_*}\,\widetilde f .
\]
By \eqref{eq:graph-plateau-value},
\(\widetilde f\le2c_H\) for large \(H_*\), and the last right
side tends to zero.  This contradiction proves \(v\ge1/2\) on
the smaller box.

It remains to control derivatives on the cutoff region without
using any growth rate for the global elliptic constants.
In a still smaller box, rescale coordinates \(x=x_0+c_Hy\) and
write \(\widetilde f=c_Hg\).  The weighted graph curvature in
the rescaled chart is \(c_HD\), and its right side becomes
\begin{align*}
 c_Hp(x,Nc_Hg,v)
  &=\frac{\epsilon a\delta^2+(1-v^2)/2}{Ng}
       -\frac{c_H^2(\epsilon a+H_*)g}{N}\\
  &=\frac{\epsilon a\delta^2+(1-v^2)/2}{Ng}
       -\frac{\epsilon a\delta^2g}{N}.
\end{align*}
Here \(1\le g\le2\) for large \(H_*\), and the angle bound gives
a uniform gradient bound.  The rescaled metric and normalized
weight have uniformly bounded derivatives; the displayed
right side and all its derivatives are uniformly bounded on
these ranges.  The local higher estimates in the proof of
Lemma~\ref{lem:graph-prescription}, applied on unit boxes in
these coordinates, consequently give uniform bounds for every
derivative of \(g\).  Scaling back gives bounds of the form
\(C_jc_H^{1-j}\) for the \(j\)-th derivatives of
\(\widetilde f\).  Interior interpolation between these polynomial
bounds and the exponential bound \eqref{eq:graph-plateau-value}, on
further fixed buffered boxes, proves \eqref{eq:graph-plateau-jets}.
Indeed the usual finite-difference interpolation of a
\(j\)-th derivative between the supremum norm and a higher
derivative norm gives an exponential factor times a power of
\(c_H^{-1}\), which is smaller than every power of \(H_*^{-1}\).

We can now cut off.  On the transition region,
\(\widetilde f\asymp c_H\), its positive-order derivatives are
smaller than every power of \(H_*^{-1}\), and the derivatives
of \(\chi\) are fixed.  On \(\{\chi>0\}\),
\[
 \frac{|d(\chi\widetilde f)|^2}{\chi\widetilde f}
 \le
 2\widetilde f\,\frac{|d\chi|^2}{\chi}
       +2\chi\,\frac{|d\widetilde f|^2}{\widetilde f}
 =O(c_H).
\]
It follows that \(A^2/f\le N^2|df|^2/f\) tends to zero.
All fixed derivatives of \(f\) are \(O(c_H)\) there.
The graph formulas therefore give \(ND,N|\II|\), and
\(N|d\log L|\) tending to zero.  They also control
\(\nabla D\), so \eqref{eq:graph-identity} gives the first
limit in \eqref{eq:graph-cutoff-small}.  Finally \(F_i=0\) in
this region, \(ND^0=-\epsilon af+A^2/(2f)\), and
\(\beta f\le H_*\,O(c_H)\to0\).  This proves every limit in
\eqref{eq:graph-cutoff-small}.  To verify the extensions at the zero
set, write \(\chi=\zeta^2\), where \(\zeta\) is smooth and flat there.
On \(\{\zeta>0\}\),
\[
 \frac{|df|_h^2}{f}
 =4\widetilde f\,|d\zeta|_h^2
   +4\zeta\langle d\zeta,d\widetilde f\rangle_h
   +\zeta^2\frac{|d\widetilde f|_h^2}{\widetilde f}.
\]
The right side extends smoothly and flatly across \(\{\zeta=0\}\),
since \(\widetilde f>0\).  The same is therefore true of
\[
 \frac{A^2}{f}
 =\frac{N^2}{1+N^2|df|_h^2}\,\frac{|df|_h^2}{f}.
\]
Moreover, \(F_i\) vanishes on a neighborhood of the zero set of
\(f\).  Thus \(ND^0\) extends smoothly by zero there, and all jets
of \(f=\zeta^2\widetilde f\) vanish on its zero set.

On the uncut region the last term
\(\varpi=\epsilon a\delta^2/f\) in
\eqref{eq:graph-core-prescription} has
\begin{equation}
 \|\varpi\|_\infty
 \le\delta\sqrt{\epsilon a(\epsilon a+H_*)}\longrightarrow0,
 \qquad
 -2NU(\varpi)
   =\frac{2\epsilon a\delta^2 L A^2}{f^2}\ge0 .
 \label{eq:graph-floor-derivative}
\end{equation}
Since \(ND^0-\beta f\ge-(\epsilon a+\beta)f\), its contribution
inside the curvature square satisfies
\begin{equation}
 (ND^0-\beta f+\varpi)^2
 \ge (ND^0-\beta f)^2
        -2\epsilon a\delta^2(\epsilon a+H_*).
 \label{eq:graph-floor-square}
\end{equation}
The last loss tends to zero.  These estimates on the core and
\eqref{eq:graph-cutoff-small} on the transition prove the
claimed bound for \(\mathrm{err}\).  The curvature estimates
use the derivative sign in \eqref{eq:graph-floor-derivative}
and the full transition estimates, as well as that bound on
the value of \(\mathrm{err}\).

Every construction has uniform bounds for fixed \(H_*\).
For Borel data, the distance to the open nonzero set of \(F_i\)
is Borel: using the enumerated charts it is the infimum of
lengths of a countable dense collection of chart paths ending
where \(F_i\ne0\).  The normalized smoothing kernels and all
subsequent scalar operations are Borel and equivariant.
Lemma~\ref{lem:graph-prescription} supplies these same
properties for the complete solution.  This proves the final
assertions.
\end{proof}

\subsection{A sweep of the root probabilities}

We now apply the localized graph to each coordinate of a root
probability map.  The order of the operations matters: a graph
enlarges the metric available to all subsequent coordinates.  We
use finitely many colors so that graphs in the same color have
disjoint supports and can be performed simultaneously.
For a root probability map, its energy is
\[
 e_h(F)=|dF|_{h,\HS}^2=\sum_i|dF_i|_h^2.
\]

Fix a smooth nonnegative nondecreasing function
\(\Psi:[0,\infty)\to[0,\infty)\) such that
\begin{equation}
 \Psi(s)=0\quad(s\le8),\qquad
 \Psi(s)\le Cs^2,\qquad |\Psi'(s)|^2\le C\Psi(s),\qquad
 \Psi(s)\ge cs^2\quad(s\ge R_\Psi),
 \label{eq:sweep-boost-profile}
\end{equation}
where \(R_\Psi\) and the constants are numerical.  Such a
function can be obtained as the square of a smooth
nondecreasing function that is zero up to \(8\), has bounded
derivative, and equals \(s-8\) from a fixed larger value on.
We use either \(\psi=0\), or \(\psi=\kappa\Psi\) for a small
constant \(\kappa>0\).  In the latter case,
\[
 \psi\le C\kappa s^2,\qquad
 |\psi'|^2\le C\kappa\psi,\qquad
 \psi(s)\ge c\kappa s^2\quad(s\ge R_\Psi).
\]

\paragraph{Notation for an ordered sweep.}
We introduce the pointwise quantities for any candidate sequence of
graphs before asserting that the sequence can be constructed.  The
pass parameters are
\begin{equation}
 \epsilon>0,\qquad N^2=l\epsilon,\qquad
 a=\tfrac12+\theta,\qquad \lambda=\tfrac12-2\theta,\qquad
 0<l\le1,\quad l\le\theta\le\tfrac1{16},
 \label{eq:sweep-parameters}
\end{equation}
with \(N>0\).  On one leaf, start with \((h_0,\phi_0,F)\), where
\(F_i\ge0\) and \(\sum_iF_i^2=1\).  Suppose nonnegative smooth
bumps \(f_i\) form a locally finite family ordered in finitely many
colors, with disjoint supports within each color, and that
\(F_i\ne0\) implies \(f_i>0\).  Attach also nonnegative functions
\(\beta_i\).  At this stage these are candidate data, with no graph
equation imposed.

Apply in color order the graph of height \(N\) times the sum of the
bumps in that color.  On the support of \(f_i\), write
\(h_{i-1},\phi_{i-1}\) for the pair before its color and
\(h_i,\phi_i\) for the pair afterward; write \(j<i\) for an earlier
color.  Let \(L_i,U_i,D_i,\II_i,H_i\) be the individual graph
quantities for \(f_i\), and let \(h_+,\phi_+\) be the final pair.
The subscript \(p\) denotes the current pass in \(W_p\) and in the
reserve \(\Pi_p\) below.
The graph identities give the output and its normalization:
\begin{equation}
 \begin{gathered}
 h_+=h_0+N^2\sum_i df_i\otimes df_i,\qquad
 W_p:=\sum_i\log L_i=\log(dV_{h_+}/dV_{h_0})\ge0,\\
 \phi_+=\phi_0-W_p,\qquad
 e^{\phi_+}dV_{h_+}=e^{\phi_0}dV_{h_0},\qquad
 r:=\sum_i f_i^2>0,\qquad Y_i=f_i/\sqrt r .
 \end{gathered}
 \label{eq:sweep-output}
\end{equation}
The sums are pointwise finite, with at most one nonzero bump per
color.  The implication \(F_i\ne0\Rightarrow f_i>0\) ensures
\(r>0\).  In all the quantities below, indices are restricted to
\(f_i>0\).

For each such index, use the longitudinal notation of the
one-coordinate calculation, with \(h=h_{i-1}\),
\(h'=h_i\), and \(h_{\mathrm{in}}=h_0\):
\[
 A_i=|U_i|_{h_{i-1}},\qquad U_i=A_i e_i,\qquad
 x_i=-\II_i(\tau_i,\tau_i),\qquad X_i=Nx_i,\qquad
 t_i=A_i^2/f_i.
\]
Here \(\tau_i=(L_i^{-1}e_i,A_i)\) and \(x_i=e_i(A_i)\) when
\(A_i>0\).  At \(A_i=0\) use any unit direction as in
\eqref{eq:graph-longitudinal}.  In the graph metric \(h_i\), define
\begin{equation}
 \begin{gathered}
 \II_i^{\mathrm{rem}}
      :=\II_i+x_i\,\tau_i^\flat\otimes\tau_i^\flat,\\
 |\II_i|^2=x_i^2+|\II_i^{\mathrm{rem}}|^2,\qquad
 \tr\II_i^{\mathrm{rem}}=H_i+x_i .
 \end{gathered}
 \label{eq:sweep-longitudinal}
\end{equation}
In particular, at \(A_i=0\) the combination
\((X_i-t_i)^2+N^2|\II_i^{\mathrm{rem}}|^2=N^2|\II_i|^2\)
is independent of the chosen direction.  The losses from comparing
the current graph with the input and final metrics are
\[
 B_i=1-|e_i|_{h_0}^2,\qquad
 C_i=
 \begin{cases}
 1-N^2|df_i|_{h_+}^2/A_i^2,&A_i>0,\\
 0,&A_i=0.
 \end{cases}
\]
Metric monotonicity and the inverse graph formula give
\(0\le B_i,C_i\le1\).

The normalization gives probability weights \(\nu_i=f_i^2/r\).
Write \(\E_\nu b_i=\sum_i\nu_i b_i\).  The mean and centered
quantities used in the estimate are
\begin{equation}
 \begin{gathered}
 m=\sum_iA_i^2,\qquad
 u_i=\frac{A_i^2}{\epsilon\nu_i},\qquad
 \bar u=\E_\nu u_i=\frac m\epsilon,\\
 \sigma_i=\frac{F_i^2}{\nu_i},\qquad \E_\nu\sigma_i=1,\\
 M_0=a(1-r)-\bar u/2,\qquad
 V_i=a(\sigma_i-1)+\lambda(u_i-\bar u),\qquad
 \Var_\nu u=\E_\nu(u_i-\bar u)^2 .
 \end{gathered}
 \label{eq:sweep-statistics}
\end{equation}
The identity for \(\E_\nu\sigma_i\) uses
\(F_i\ne0\Rightarrow f_i>0\); the individual ratios \(\sigma_i\)
need not be at most one.

Finally, define two nonnegative quantities:
\begin{align}
 T_0={}&\sum_i A_i^2\left[
 C_i+4a^2\sigma_iB_i+(L_i-1)\{2a(\sigma_i+r)+u_i\}\right],
 \label{eq:sweep-T}\\
 \mathcal X={}&\sum_i\left[
 (X_i-t_i)^2+N^2|\II_i^{\mathrm{rem}}|^2
       +N^2|d\log L_i|_{h_i}^2+\beta_i^2f_i^2\right].
 \label{eq:sweep-X}
\end{align}
The quantity \(T_0\) records the gains from the metric comparisons
and from each graph's stretching.  Its full coefficient will yield
the sharp quadratic term \(m^2/n\) in Section~\ref{sec:entropy}.
The quantity \(\mathcal X\) retains the derivative squares and the
applied coefficients.  Together with the mean and variance squares
below, it controls graph curvature; those moment squares also control
the normalization.

\begin{proposition}[One sweep]
\label{prop:sweep}
There is a numerical \(c>0\) with the following property.
For a sweep without boost, set \(\psi=0\).  For a boosted sweep,
given \(\epsilon_*>0\), there are constants
\(\kappa>0\), \(c_*>0\), and \(C_*<\infty\), depending only on
\(\epsilon_*\), for which one may use \(\psi=\kappa\Psi\) whenever
\(\epsilon\ge\epsilon_*\).

Let \(h_0,\phi_0,F\) be a smooth equivariant Borel family on
\(\widetilde M\times\Omega\), with bounds uniform in the lifted
charts for the instance.  Suppose \(h_0\) is complete, its
geometry and the positive-order derivatives of \(\phi_0\) have
the uniform bounds above, and
\[
 F=(F_i)_{i\in\mathcal L},\qquad F_i\ge0,\qquad
 \sum_iF_i^2=1
\]
has finitely many free label orbits and uniformly locally finite
compact coordinate supports.  For each orbit representative,
suppose all its transverse supports lie in a fixed compact
subset of \(\widetilde M\).  Let \(h_{\mathrm{ref}}\le h_0\)
be a complete equivariant Borel leafwise metric whose coefficients,
inverse, and all coordinate derivatives have uniform bounds in the
same lifted buffered charts as \(h_0\).  Suppose a Borel
equivariant coloring
\(\chi:\mathcal L\times\Omega\to\{1,\ldots,k\}\), \(k<\infty\),
has been chosen such that, for each retained \(\omega\) and distinct
labels \(i,j\) with \(\chi(i,\omega)=\chi(j,\omega)\), the
\(R_{\mathrm{loc}}\)-neighborhoods in \(h_{\mathrm{ref}}(\omega)\)
of \(\supp F_i(\cdot,\omega)\) and \(\supp F_j(\cdot,\omega)\)
are disjoint.  Fix the pass parameters in
\eqref{eq:sweep-parameters}, with \(\epsilon\ge\epsilon_*\) in
the boosted case.

For any \(\delta_0,\eta_0>0\), there is a candidate sweep as above
with smooth nonnegative bumps \(f_i\le2\), supported in the
indicated enlargements, positive on the supports of \(F_i\), and
flat on their zero sets.  Its output is given by
\eqref{eq:sweep-output}.  The lower bound for \(r\) is uniform
for the instance.  The bumps, coefficients, all intermediate and
final graph metrics and weights, their smooth graph quantities, and
\(Y\) are equivariant and Borel and retain the corresponding uniform
chart bounds, with the positive-order convention for weights.
The map \(Y\) is again a
root probability map with the same label orbits.

For each \(i\), the coefficient furnished by
Lemma~\ref{lem:graph-localization} has the form
\begin{equation}
 \beta_i=\beta_{s,i}+\epsilon\psi(R_i),\qquad
 R_i=\sum_{j<i}f_j^2 ,
 \label{eq:sweep-beta}
\end{equation}
Here \(\beta_{s,i}\) vanishes near \(\supp F_i\), and the coefficients
satisfy \eqref{eq:graph-spatial-cutoff} in \(h_{i-1}\).
Writing
\[
 ND_i^0=\epsilon a(F_i^2/f_i-f_i)+A_i^2/(2f_i)
\]
on \(\{f_i>0\}\), extended by zero off the bump, one has
\begin{equation}
 ND_i=ND_i^0-\beta_i f_i+\mathrm{err}_i,\qquad
 \sum_i|\mathrm{err}_i|^2\le\eta_0 .
 \label{eq:sweep-prescription-error}
\end{equation}

Define the nonnegative reserve
\begin{equation}
 \begin{split}
 \Pi_p:={}&
 T_0+c\epsilon\bigl(M_0^2+\E_\nu V_i^2+\theta\Var_\nu u\bigr)
       \\
 &+c\frac r\epsilon\mathcal X
       +c_*\epsilon r^6\indic_{\{r\ge C_*\}},
 \end{split}
 \label{eq:sweep-positive-reserve}
\end{equation}
where the final term is omitted in a sweep without boost.  Then,
pointwise,
\begin{align}
 Q:=\frac{rN^2}{\epsilon}\bigl[
 &\wscal(h_+,\phi_+)-\wscal(h_0,\phi_0)
       +\epsilon e_{h_0}(F)-\epsilon e_{h_+}(Y)\bigr]
 \notag\\
 &\ge\Pi_p-4\theta^2m-\frac r\epsilon\delta_0 .
 \label{eq:sweep-inequality}
\end{align}
The numerical constant \(c\) is independent of the instance, the
number of colors, and the pass parameters.  The errors
\(\delta_0,\eta_0\) can be prescribed independently before the pass.
\end{proposition}

\begin{proof}
Process the colors in their prescribed order.  At a label in
the current color, the sum \(R_i\) of earlier bumps and all its
derivatives have common bounds for that color, so the function
\(\alpha_i=\epsilon\psi(R_i)\) is fixed and uniformly smooth.
Apply Lemma~\ref{lem:graph-localization} with this
\(\alpha_i\), using a common sufficiently large \(H_*\) for
the color.  Uniqueness of the complete solutions and the
equivariant cutoff construction make all choices equivariant.
Their Borel dependence follows from the same lemma.  Bumps in
the same color are separated, so their sum gives exactly the
individual graph at each point of their supports.  The graph
metric and weight bounds hold in the same lifted charts as the input
bounds.  The reference metric is unchanged during the pass, so the
simultaneous chart control required by
Lemma~\ref{lem:graph-localization} persists for the next color.

At any point some \(F_i\) is nonzero, and its cutoff is one
there.  The lower barrier for that color gives
\(f_i\ge c_{\min}>0\).  Taking the minimum over the finitely
many colors gives the stated instance-uniform lower bound for
\(r\); also \(r\le4k\).  Thus \(Y\) is smooth with uniform
bounds and \(\sum_iY_i^2=1\).  The graph determinant and
density identities prove \eqref{eq:sweep-output}.
By making the error in each color small enough that its square
is at most \(\eta_0/k\), the last conclusion of
Lemma~\ref{lem:graph-localization} gives
\eqref{eq:sweep-prescription-error}.

We prove the curvature estimate in the units before
multiplication by \(r/\epsilon\).  Write
\[
 G=N^2\bigl[\wscal(h_+,\phi_+)-\wscal(h_0,\phi_0)
       +\epsilon e_{h_0}(F)-\epsilon e_{h_+}(Y)\bigr],
 \qquad Z_i=ND_i^0,
\]
so \(Q=(r/\epsilon)G\).  For each \(i\), the one-coordinate calculation
applies with \(h=h_{i-1}\), \(h'=h_i\), and
\(h_{\mathrm{in}}=h_0\).  We first account for the difference between
the ideal target \(Z_i\) for \(ND_i\) and the realized prescription.

For the moment omit the floor in
\eqref{eq:graph-core-prescription}.  Replacing \(Z_i\) by
\(Z_i-\beta_i f_i\) in the first two terms of
\eqref{eq:graph-identity}, after multiplying by \(N^2\),
adds exactly
\begin{align*}
 &(Z_i-\beta_i f_i)^2-2NU_i(Z_i-\beta_i f_i)
                 -Z_i^2+2NU_i(Z_i)\\
 &\hspace{1cm}=\beta_i^2f_i^2
       +2\epsilon a\beta_i(f_i^2-F_i^2)
       +(2L_i-1)\beta_iA_i^2
       +2NA_if_i e_i(\beta_i).
\end{align*}
Since \(L_i\ge1\), this is at least
\begin{equation}
 \mathcal B_i:=
 \beta_i^2f_i^2+2\epsilon a\beta_i f_i^2+\beta_i A_i^2
       -2\epsilon a\beta_iF_i^2
       -2NA_if_i|d\beta_i|_{h_{i-1}}.
 \label{eq:sweep-beta-correction}
\end{equation}

On the uncut region the actual floor has nonnegative
differentiated contribution by
\eqref{eq:graph-floor-derivative}, and its square loses at
most the quantity in \eqref{eq:graph-floor-square}.
On a cutoff transition \(F_i=0\) and \(d\beta_i=0\).
There \eqref{eq:graph-cutoff-small} makes the actual scalar
change small.  It also makes every term on the right of the
single-step expansion just obtained small:
\(\beta_i A_i^2=(\beta_i f_i)t_i\),
the longitudinal terms are controlled by \(N|\II_i|\) and
\(t_i\), and the other terms are controlled by
\(Z_i,\beta_i f_i,A_i\).  The first term in
\eqref{eq:graph-ideal-derivative} is zero.  Thus this
expanded lower bound still holds on a cutoff transition with
an arbitrarily small additive error in these units.  It is
also valid off the bumps, where all its terms vanish.
Choose the floor and transition errors for each color so their
sum is at most \(\delta_0\) at every point.  Telescoping
\eqref{eq:graph-identity} over the colors gives
\begin{align}
 G\ge \sum_i\bigl[&
 Z_i^2-2NU_i(Z_i)+N^2|\II_i|^2
       +N^2|d\log L_i|_{h_i}^2
       +\epsilon N^2|dF_i|_{h_0}^2+\mathcal B_i\bigr]
 \notag\\
 &-\epsilon N^2e_{h_+}(Y)-\delta_0 .
 \label{eq:sweep-raw}
\end{align}
This argument uses the full differentiated floor and cutoff
estimates, rather than differentiating the value error in
\eqref{eq:sweep-prescription-error}.

The input energy is already included in
\eqref{eq:graph-ideal-bound}.  For the output, normalization is an
orthogonal projection in the target:
\[
 dY=r^{-1/2}(I-Y\otimes Y)\,df .
\]
The rank-one inverse identity and metric monotonicity give
\[
 N^2|df_i|_{h_i}^2=A_i^2,\qquad
 N^2|df_i|_{h_+}^2=A_i^2(1-C_i)\le A_i^2 .
\]
Consequently
\begin{equation}
 N^2e_{h_+}(Y)
       \le\frac1r\sum_i A_i^2(1-C_i).
 \label{eq:sweep-output-energy}
\end{equation}

Put
\[
 K_i=(X_i-t_i)^2+N^2|\II_i^{\mathrm{rem}}|^2
                 +N^2|d\log L_i|_{h_i}^2,\qquad
 g_i=a(\sigma_i-r)+u_i/2 .
\]
The definitions imply
\[
 Z_i=\frac{\epsilon f_i}{r}g_i,\quad
 \frac r\epsilon\sum_iZ_i^2=\epsilon\E_\nu g_i^2,\quad
 \frac r\epsilon(L_i-1)t_i^2=(L_i-1)A_i^2u_i .
\]
Use these identities, \eqref{eq:graph-ideal-bound}, and
\eqref{eq:sweep-output-energy} in \eqref{eq:sweep-raw}.
After multiplication by \(r/\epsilon\), the result is
\begin{align}
 Q\ge{}&T_0+\frac r\epsilon\sum_iK_i
       +\epsilon\E_\nu g_i^2
       +\sum_i A_i^2\bigl[(2a-4a^2)\sigma_i+2ar-1\bigr]
       +\frac r\epsilon\left(\sum_i\mathcal B_i-\delta_0\right).
 \label{eq:sweep-expanded}
\end{align}
Indeed \(rF_i^2/f_i^2=\sigma_i\).  The terms with \(B_i\) and
\(L_i-1\) in \eqref{eq:graph-ideal-bound} become the corresponding
terms of \(T_0\), while the output contributes
\(-A_i^2+A_i^2C_i\).  The remaining height terms give the displayed
quadratic expression.

The remaining quadratic expression has an exact mean and
variance decomposition.  Use \(\E_\nu\sigma_i=1\) and
\(\E_\nu u_i=\bar u\) from \eqref{eq:sweep-statistics}, and write
\(w_i=u_i-\bar u\).  Dividing the quadratic expression in
\eqref{eq:sweep-expanded} by \(\epsilon\), its contribution
from the means is
\begin{align*}
 [a(1-r)+\bar u/2]^2+\bar u(2a-4a^2+2ar-1)
 &= [a(1-r)-\bar u/2]^2-(2a-1)^2\bar u\\
 &=M_0^2-4\theta^2\bar u .
\end{align*}
The coefficient of
\(\E_\nu[w_i(\sigma_i-1)]\) in the centered expression is
\[
 a+(2a-4a^2)=3a-4a^2=2a\lambda .
\]
Completing its square therefore gives the identity
\begin{align}
 \epsilon\E_\nu g_i^2
 &+\sum_i A_i^2\bigl[(2a-4a^2)\sigma_i+2ar-1\bigr] \notag\\
 &=
 \epsilon\left(M_0^2+\E_\nu V_i^2+
                   (\tfrac14-\lambda^2)\Var_\nu u\right)
        -4\theta^2m .
 \label{eq:sweep-mean-identity}
\end{align}
Here
\(\tfrac14-\lambda^2=2\theta-4\theta^2\ge7\theta/4\).
This supplies a reserve for all three squares in the
proposition.  We finish by paying the terms in
\(\mathcal B_i\) from fixed fractions of this reserve and of
the nonnegative \(\beta_i\) terms.

Write \(s_i=\beta_{s,i}\) and \(b_i=\epsilon\psi(R_i)\), so
\(\beta_i=s_i+b_i\) with both summands nonnegative.
In particular \(\beta_i^2\ge s_i^2+b_i^2\).
The negative \(F_i^2\) term in
\eqref{eq:sweep-beta-correction} has no \(s_i\) part, because
\(s_i\) vanishes near the support of \(F_i\).
The spatial derivative term satisfies
\begin{align}
 2NA_if_i|ds_i|
 &\le2c_0 A_if_i(\theta\epsilon^{3/2}+s_i^{3/2}) \notag\\
 &\le c_0(s_i^2f_i^2+s_iA_i^2)
       +Cc_0\theta(\epsilon^2f_i^2+A_i^4/f_i^2).
 \label{eq:sweep-spatial-cost}
\end{align}
For the first pair of terms use \(2xy\le x^2+y^2\).
For the other pair, divide by \(\epsilon^2f_i^2\) and use
\(t\le1+t^4\) with \(t=A_i/(\sqrt\epsilon f_i)\).
This derivative cost occurs only in rows with \(F_i=0\), hence
\(\sigma_i=0\).  In those rows,
\[
 V_i=-a+\lambda w_i,\qquad
 ar+\bar u/2=a-M_0=\lambda w_i-V_i-M_0 .
\]
Since \(a\ge1/2\) and \(|\lambda|\le1/2\), these equations give
\[
 r+u_i\le2(|M_0|+|V_i|+|w_i|),\qquad
 r^2+u_i^2\le12(M_0^2+V_i^2+w_i^2).
\]
After multiplication by \(r/\epsilon\), the last cost in
\eqref{eq:sweep-spatial-cost}, summed over these rows, is
\[
 Cc_0\epsilon\theta
     \sum_{\{\sigma_i=0\}}\nu_i(r^2+u_i^2).
\]
The preceding bound shows that it is a small fraction of the
three squares in \eqref{eq:sweep-mean-identity} when \(c_0\)
is a sufficiently small numerical constant.  The first cost
in \eqref{eq:sweep-spatial-cost} is paid by a small fraction
of \(\beta_i^2f_i^2+\beta_iA_i^2\).
This proves the required absorption when \(\psi=0\), leaving
a fixed fraction of every asserted square.

Suppose now that \(\psi=\kappa\Psi\) and
\(\epsilon\ge\epsilon_*\).  Metric addition over earlier
colors implies, as quadratic forms,
\[
 h_{i-1}\ge N^2\sum_{j<i}df_j\otimes df_j .
\]
For a unit vector in \(h_{i-1}\), the vector
\((Ndf_j)_{j<i}\) consequently has Euclidean norm at most one.
Cauchy--Schwarz applied to
\(dR_i=2\sum_{j<i}f_jdf_j\) proves
\begin{equation}
 |dR_i|_{h_{i-1}}\le2\sqrt{R_i}/N .
 \label{eq:sweep-prefix-gradient}
\end{equation}
For any fixed small \(\zeta>0\), the derivative bound in
\eqref{eq:sweep-boost-profile} and Young's inequality give
\[
 2NA_if_i|db_i|
 \le4\epsilon A_if_i\sqrt{R_i}\,|\psi'(R_i)|
 \le\zeta\epsilon\psi(R_i)A_i^2
          +C_\zeta\epsilon\kappa f_i^2R_i .
\]
The first term is a small fraction of \(\beta_iA_i^2\).
For the other costs, use
\[
 \sum_i f_i^2R_i
       =\tfrac12\left(r^2-\sum_i f_i^4\right)\le r^2/2,
 \qquad
 \sum_i\psi(R_i)F_i^2\le C\kappa r^2\sum_iF_i^2=C\kappa r^2.
\]
The profile and its derivative vanish up to \(8\).  Thus the
sum of all remaining boost costs in the units of \(G\) is at
most
\begin{equation}
 C(\epsilon+\epsilon^2)\kappa r^2\indic_{\{r\ge8\}}.
 \label{eq:sweep-boost-cost}
\end{equation}

Two positive terms cover this cost.  First, for \(r\ge8\),
\[
 -M_0=a(r-1)+\bar u/2\ge7r/16,
 \qquad
 \frac{\epsilon^2}{r}M_0^2\ge c\epsilon^2r .
\]
The latter is the mean square in
\eqref{eq:sweep-mean-identity} expressed in the units of \(G\).
Second, the increments \(f_i^2\) in \(R_i\) are at most \(4\).
For \(r\) beyond a fixed threshold, the rows with
\(R_i\ge r/2\) carry total increment at least
\(r/2-4\ge r/4\).  In fact the first prefix crossing \(r/2\)
overshoots it by at most \(4\).  Once \(r/2\ge R_\Psi\), the
lower profile bound on each of these rows gives
\begin{equation}
 \sum_i f_i^2 b_i^2
   =\epsilon^2\sum_i f_i^2\psi(R_i)^2
       \ge c\epsilon^2\kappa^2 r^5
       \qquad(r\ge C_*).
 \label{eq:sweep-boost-square}
\end{equation}

To make the parameter choice explicit, let
\(K_*=C(1+\epsilon_*^{-1})\).  The cost
\eqref{eq:sweep-boost-cost} is at most
\(K_*\epsilon^2\kappa r^2\) on \(r\ge8\).
Choose \(\kappa\) so small that
\(\kappa^{-1/2}\ge C_*\) and \(K_*\sqrt\kappa\) is smaller
than the fixed fractions of the preceding two reserves
available for this purpose.  If
\(8\le r\le\kappa^{-1/2}\), this cost is at most
\(K_*\sqrt\kappa\,\epsilon^2r\).  If
\(r\ge\kappa^{-1/2}\), it is at most
\(K_*\sqrt\kappa\,\epsilon^2\kappa^2r^5\).
It is therefore covered in both ranges.  Fixed fractions of
the mean square and of \eqref{eq:sweep-boost-square} remain.
One can simultaneously retain a fixed fraction of
\(\sum_i\beta_i^2f_i^2\) for \(\mathcal X\): reserve, for
example, one half of this sum at the outset and use
\(\beta_i^2\ge s_i^2+b_i^2\) on the other half for the
spatial and boost estimates.  Take \(c_0,\zeta,\kappa\) small
enough for these fixed allocations.

Insert the absorptions into \eqref{eq:sweep-expanded} and
\eqref{eq:sweep-mean-identity}.  A fixed fraction of
\(\sum_iK_i+\sum_i\beta_i^2f_i^2=\mathcal X\) and of the
three mean and variance squares survives.  The remaining
part of \eqref{eq:sweep-boost-square}, after multiplication
by \(r/\epsilon\), is
\(c_*\epsilon r^6\indic_{\{r\ge C_*\}}\), with
\(c_*\) a positive multiple of \(\kappa^2\).
The terms \(T_0\) and \(-4\theta^2m\) have not been spent.
Choose the constant \(c\) in \eqref{eq:sweep-positive-reserve}
below the fixed numerical fractions that survive for \(\mathcal X\)
and the mean and variance squares.  These fractions are independent
of \(\epsilon_*\), because \(\kappa\) was chosen to fit the fixed
allocations.
This proves \eqref{eq:sweep-inequality}.

Finally, the choices of errors impose no estimate on the
growth of analytic constants from one color to the next.
For a fixed color, the reference buffers are fixed for the
pass, and the smaller analytic boxes and all previous
geometry bounds are fixed before its \(H_*\) is chosen.
There are finitely many colors, so the raw scalar error can
be made at most \(\delta_0/k\) per color and the squared
value error at most \(\eta_0/k\) per color.  Later choices
of \(H_*\) may be arbitrarily larger.  The summed errors
are therefore precisely those asserted in
\eqref{eq:sweep-prescription-error} and
\eqref{eq:sweep-inequality}.
\end{proof}
